% Part: sets-functions-relations
% Chapter: relations
% Section: operations

\documentclass[../../../include/open-logic-section]{subfiles}

\begin{document}

\olfileid[fa-IR]{sfr}{rel}{ops}
\olsection{عملیات روی روابط}

اغلب مفید است که روابط را تغییر دهیم یا با هم ترکیب کنیم. در
\olref[sfr][rel][ord]{prop:stricttopartial}، \emph{اجتماعِ}
روابط را بررسی کردیم، که صرفاً اجتماع دو رابطه است وقتی آن دو را
به‌منزلهٔ مجموعه‌هایی از زوج‌ها در نظر بگیریم. به همین ترتیب، در
\olref[sfr][rel][ord]{prop:partialtostrict}، تفاضل نسبی روابط را
بررسی کردیم. در اینجا چند عمل دیگر آمده است که می‌توانیم روی روابط انجام دهیم.

\begin{defn}\ollabel{relationoperations}
فرض کنید $R$ و $S$ رابطه باشند و $A$ مجموعه‌ای دلخواه باشد.

\emph{وارون} $R$ برابر است با $R^{-1} = \Setabs{\tuple{y, x}}{\tuple{x,
    y} \in R}$.

\emph{حاصل‌ضرب نسبی} $R$ و $S$ برابر است با $(R \mid S) =
\{\tuple{x, z} : \exists y(Rxy \land Syz)\}$.

\emph{تحدید} $R$ به $A$ برابر است با $\funrestrictionto{R}{A}= R
\cap A^2$.

\emph{اعمال} $R$ بر $A$ برابر است با $\funimage{R}{A} = \{y :
(\exists x \in A)Rxy\}$
\end{defn}

\begin{ex}
فرض کنید $S \subseteq \Int^2$ رابطهٔ جانشینی روی~$\Int$ باشد؛ یعنی
$S = \Setabs{\tuple{x, y} \in \Int^2}{x + 1 = y}$، به‌طوری‌که $Sxy$ اگر و تنها اگر $x + 1 = y$.

$S^{-1}$ رابطهٔ پیشینی روی $\Int$ است؛ یعنی
$\Setabs{\tuple{x,y}\in\Int^2}{x -1 =y}$.

$S\mid S$ برابر است با
$ \Setabs{\tuple{x,y}\in\Int^2}{x + 2 =y}$

$\funrestrictionto{S}{\Nat}$ رابطهٔ جانشینی روی~$\Nat$ است.

$\funimage{S}{\{1,2,3\}}$ برابر است با $\{2, 3, 4\}$.
\end{ex}

\begin{defn}[بستار تراگذری]فرض کنید $R \subseteq A^2$ یک رابطهٔ دوتایی باشد.

\emph{بستار تراگذریِ}~$R$ برابر است با $R^+ = \bigcup_{0 < n \in
\Nat} R^n$، که در آن به‌صورت بازگشتی $R^1 = R$ و $R^{n+1} = R^n
\mid R$ را تعریف می‌کنیم.

\emph{بستار بازتابی و تراگذریِ} $R$ برابر است با $R^* = R^+ \cup
\Id{A}$.
\end{defn}

\begin{ex}
رابطهٔ جانشینی $S \subseteq \Int^2$ را در نظر بگیرید. $S^2xy$ اگر و تنها اگر $x + 2 =
y$، و $S^3xy$ اگر و تنها اگر $x + 3 = y$، و به همین ترتیب. پس $S^+xy$ اگر و تنها اگر $x + n = y$ برای یک
$n \geq 1$ برقرار باشد. به عبارت دیگر، $S^+xy$ اگر و تنها اگر $x < y$، و $S^*xy$ اگر و تنها اگر $x \le
y$.
\end{ex}

\begin{prob}
نشان دهید که بستار تراگذریِ $R$ در واقع تراگذری است.
\end{prob}

\end{document}
